\songtitle{Vous ?! — Ça te dérange ?}

\setin{2005-06-21}
\tag*{2026}\tag{sophisti-pop}\tag{french-lyrics}\tag{ricochets}\tag{fabian-muller}\tag{sophie-durand}
% -> ricochets.tex
\begin{notes}
Lyrics by Claude based on \emph{Ricochets} by Carlos Thompson -- chapter 1 (\emph{The Hangover}): Müller meets Sophie Durand at the hotel bar under the Old King Cole mural in New York; the slip from \emph{vous} to \emph{tu}, the contract talk she finds boring, and the walk along 55th Street for a New York pizza. Retrieved from rataflechera's Suno page (V6, 15 November 2026).
\end{notes}
\begin{style}
2000s sophistipop with chanson française and cabaret elements, light bossa swing at 98 BPM; glossy late-night mix; brushed drums, upright bass, Rhodes, muted trumpet, nylon guitar, accordion touches and string pads; warm male baritone with a light Spanish accent and sultry female alto with a southern French lilt, clipped rhythmic staccato phrasing, layered octave doubles and playful call-and-response.
\end{style}
\begin{lyrics}
\songpart{Intro}
\annotation{Rhodes and muted trumpet; glasses, a bar murmur}

\songpart{Verse 1}
\annotation{Fabián - male:}\\
Bonjour, mademoiselle Durand, je vous attendais,\\
Monsieur Müller, à votre entière disposition.\\
Appelez-moi Fabián, si cela vous plaît,\\
on laisse au vestiaire les formules et les façons.
\annotation{Sophie - female:}\\
Enchantée, Fabián. Appelez-moi Sophie.\\
Et Monsieur Michel ? Ne l'attendez pas ce soir.
\annotation{Fabián - male:}\\
Il ne vient pas ?
\annotation{Sophie - female:}\\
Non. Vous m'avez pour vous seul, voici.

\songpart{Verse 2}
\annotation{Fabián - male:}\\
Sous le vieux roi Cole qui rit sur le mur,\\
une robe lie-de-vin, des talons de sept pouces,\\
vous me dépassez d'une tête, c'est sûr,\\
et l'accent du Midi rend la langue plus douce.
\annotation{Sophie - female:}\\
Et vous, Monsieur, sous le français bien mis,\\
j'entends l'Antioquia, les montagnes, les cafés.
\annotation{Fabián - male:}\\
Ça se voit tant ?
\annotation{Sophie - female:}\\
Ça s'entend, mon ami.

\songpart{Pre-Chorus}
\annotation{both:}\\
Un premier verre, on se vouvoie,\\
un deuxième verre, on ne sait pas,\\
quelque part entre deux glaçons\\
on a perdu les « vous » et les façons.

\songpart{Chorus}
\annotation{Fabián - male, spoken over the band:}\\
It's time to talk business, Sophie... ça vous dérange ?
\annotation{Sophie- female, faking offence:}\\
Vous ?!
\annotation{Fabián - male:}\\
Ça te dérange ?
\annotation{both:}\\
Une lettre de moins, un monde qui change,\\
entre le vous et le tu, un pas de danse étrange.
\annotation{Sophie - female:}\\
Vous ?!
\annotation{Fabián - male:}\\
Pardon. Ça te dérange ?

\songpart{Verse 3}
\annotation{Fabián - male, in English, crisp:}\\
Two shows in Medellín for Colombiamoda,\\
one in Bogotá, where I'll look after you myself;\\
Leandro Restrepo handles all the logistics,\\
Monsieur Michel's conditions, list by list...
\annotation{Sophie - female, a fake yawn:}\\
Boring.\\
C'est pour ça que j'ai un manager.\\
Moi, je veux m'amuser, je veux de l'air,\\
et une pizza, la fameuse, de New York.
\annotation{Fabián - male:}\\
Une pizza ? Alors, une pizza. D'accord.

\songpart{Chorus}
\annotation{Sophie - female:}\\
Monsieur Müller...
\annotation{Fabián - male, mock-disappointed:}\\
Monsieur ?
\annotation{Sophie - female, laughing:}\\
Ça te dérange ?
\annotation{both:}\\
Une lettre de moins, un monde qui change,\\
entre le vous et le tu, un pas de danse étrange.
\annotation{Fabián - male:}\\
Tu me tutoies ?
\annotation{Sophie - female:}\\
Ça te dérange ?

\songpart{Bridge}
\annotation{strings rise, accordion enters}
\annotation{Fabián - male:}\\
Premier soir d'été, la chaleur sur l'asphalte,\\
mon costume de soie fait comme si c'était Miami.
\annotation{Sophie - female:}\\
On marche sur la Cinquante-Cinquième, pas de halte,\\
vers Broadway, la ville est à nous, à demi.
\annotation{Fabián - male:}\\
Et l'anglais reste au bar avec les contrats.
\annotation{Sophie - female:}\\
Dans la rue on reparle français, toi et moi.

\songpart{Final Chorus}
\annotation{both:}\\
Vous ? Tu ? Qu'importe, la nuit nous arrange,\\
une lettre de moins, un monde qui change.
\annotation{Sophie - female:}\\
Je te tutoie...
\annotation{Fabián - male:}\\
Ça me dérange ?
\annotation{both, softly:}\\
Non. Pas du tout. Pas du tout.

\songpart{Outro}
\annotation{muted trumpet, taxi horns far away}
\annotation{Sophie - female, spoken:}\\
Une pizza, Fabián.
\annotation{Fabián - male, spoken:}\\
Une pizza sera.
\annotation{End}
\end{lyrics}

